%!TEX program = xelatex
\documentclass[12pt, pdf]{style} 
%Wähle oben entweder pdf oder print, default ist print. Pdf generiert einseite Dokumente, die an Bildschirmen besser aussehen, print ist für zweiseitigen druck (Buchdruck) und lässt daher Kapitel immer auf der rechten Seite beginnen. Das kann in PDFs nervige Leerseiten erzeugen, daher kann man das hier einstellen.
\tikzset{external/prefix=tikz/Main/}
\begin{otherlanguage}{ngerman}
\iffalse ANMERKUNGEN

- die Vorlage muss folgendermaßen kompiliert werden da sonst unter Umständen die Abkürzungen oder das Literaturverzeichnis nicht funktionieren:
    1. XeLaTex
    2. Makeglossaries
    3. Biber
    4. XeLaTex
    5. XeLaTex
- die Formatierung kann in style.cls editiert werden
- Referenzen im BibTeX Format können z. B. über Google Scholar gefunden und in libraryJR.bib eingefügt werden
\fi
\end{otherlanguage}

\graphicspath{{4_images/}}
\addbibresource{3_backmatter/library/libraryJR.bib}

\newcommand{\FirstName}{Vorname}
\newcommand{\LastName}{Nachname}
\newcommand{\Matriculation}{Matrikelnummer}
\newcommand{\Place}{Ort}
\newcommand{\Signature}{\centering\includegraphics[height=1.5 cm]{Unterschrift.png}}
\newcommand{\Date}{\today}
\newcommand{\titleEN}{Titel auf Englisch}
\newcommand{\titleDE}{Titel auf Deutsch}
\newcommand{\instituteEN}{Institutname auf Englsich}
\newcommand{\instituteDE}{Institutname auf Deutsch}
\newcommand{\texttypeDE}{Art der Arbeit auf Deutsch}
\newcommand{\texttypeEN}{Art der Arbeit auf Englisch}
\newcommand{\supervisorA}{Name}
\newcommand{\supervisorB}{Name}
\newcommand{\profA}{Name}
\newcommand{\profB}{Name}


\begin{document}
% Titlepage, table of contents, etc.
    \frontmatter
    \pagenumbering{Roman}  
    \begin{titlepage}
        \begin{center}
{\Large RWTH Aachen University}\\[2cm]
{\Huge Titel der Arbeit}\\[1cm]
{\large Bachelorarbeit / Masterarbeit}\\[2cm]
Autor: Dein Name\\
Betreuer: Dr. Mustermann\\
Datum: \today
\end{center}
    \end{titlepage}
    %!TEX root = ../../main.tex

% Erste Seite hinter Titelseite
% Kurzfassung der Arbeit

\begin{center} {\LARGE Abstract} \end{center}

%Abstract hier



% present tense: Give background information; 1 general sentence about the socioeconomic or scientific relevance of your work
%present perfect: State a problem in 1 sentence with the conventional way of doing research in your particular discipline. Why is the conventional way of doing things in your field problematic? Expensive, complicated, imprecise?
%present tense: State your objective of your thesis: state the aim; This study presents / aims to
%past tense: State your methods (most important)
%past tense: State your most significant results (if you have enough space left discuss them briefly) 
%State your conclusion: take-home message for your readers; what do you want them to learn?

\cleardoublepage

    %!TEX root = ../../main.tex

% Zweite Seite des Dokuments mit Aufgabenstellung

\begin{center}
\huge{Task Description}


\end{center}
\normalsize 


\cleardoublepage

    
    %!TEX root = ../../main.tex

% Seite mit eidesstattlicher Erklärung
\begin{otherlanguage}{ngerman}

% Erkärung nach Vorlage von AVT/PT
\begin{center} {\LARGE Eidesstattliche Erklärung} \end{center}


\noindent\rule[1ex]{30ex}{1pt} \hspace{15ex} \rule[1ex]{30ex}{1pt}\vspace{-0.2cm}\\
\phantom{}\hspace{6ex} Name, Vorname  \hspace{20.5ex} Matrikelnummer (freiwillige Angabe)\vspace{0.1cm}\\
Ich versichere hiermit an Eides Statt, dass ich die vorliegende Arbeit/Bachelorarbeit/Masterarbeit* mit dem Titel \vspace{0.2cm}\\
\noindent \rule[1ex]{\textwidth}{1pt}\\
\noindent \rule[1ex]{\textwidth}{1pt}\\
\noindent \rule[1ex]{\textwidth}{1pt}\\
selbstständig und ohne unzulässige fremde Hilfe (insbes. akademisches Ghostwriting) erbracht habe.
Ich habe keine anderen als die angegebenen Quellen und Hilfsmittel benutzt; dies umfasst
insbesondere auch Software und Dienste zur Sprach-, Text- und Medienproduktion. Ich erkläre, dass
für den Fall, dass die Arbeit in unterschiedlichen Formen eingereicht wird (z.B. elektronisch, gedruckt,
geplottet, auf einem Datenträger) alle eingereichten Versionen vollständig übereinstimmen. Die Arbeit
hat in gleicher oder ähnlicher Form noch keiner Prüfungsbehörde vorgelegen.\vspace{-0.2cm}\\

\noindent\rule[1ex]{30ex}{1pt} \hspace{15ex} \rule[1ex]{30ex}{1pt}\vspace{-0.2cm}\\
\phantom{}\hspace{6ex} Ort, Datum  \hspace{37.5ex} Unterschrift \vspace{0.1cm}\\
*Nichtzutreffendes bitte streichen\vspace{0.1cm}\\
\textbf{Belehrung:}\\
\textbf{\S 156 StGB: Falsche Versicherung an Eides Statt}

\noindent Wer vor einer zur Abnahme einer Versicherung an Eides Statt zuständigen Behörde eine solche Versicherung
falsch abgibt oder unter Berufung auf eine solche Versicherung falsch aussagt, wird mit Freiheitsstrafe bis zu drei
Jahren oder mit Geldstrafe bestraft.

\textbf{\S 161 StGB: Fahrlässiger Falscheid; fahrlässige falsche Versicherung an Eides Statt}

\noindent(1) Wenn eine der in den \S\S 154 bis 156 bezeichneten Handlungen aus Fahrlässigkeit begangen worden ist, so tritt Freiheitsstrafe bis zu einem Jahr oder Geldstrafe ein.

\noindent (2) Straflosigkeit tritt ein, wenn der Täter die falsche Angabe rechtzeitig berichtigt. Die Vorschriften des § 158 Abs. 2 und 3 gelten entsprechend.

\noindent Die vorstehende Belehrung habe ich zur Kenntnis genommen:\vspace{0.6cm}\\
\noindent\rule[1ex]{30ex}{1pt} \hspace{15ex} \rule[1ex]{30ex}{1pt}\vspace{-0.2cm}\\
\phantom{}\hspace{6ex} Ort, Datum  \hspace{37.5ex} Unterschrift \vspace{0.1cm}\newpage

\begin{center} {\LARGE Declaration} \end{center}

I hereby declare to have listed all aids used in this thesis alongside their application below. This concerns in particular software and services for language, text, and media production. 
I have been made aware by my supervisor of the proper and improper use of such aids in the thesis.\newline
In this thesis, I have
\begin{itemize}
	\item[$\Box$] In particular not used software and services for language, text, and media production. 
	\item[$\Box$] 
	used the following aids (Please also list how you applied each aid):\\\\


\centering
\begin{tabular}{lll} 
\toprule
Section & Tool(s) used & Reason for use  \\
\midrule
1       & DEEPl        & Grammar/ Spell checker, Reformulating            \\
            
\bottomrule
\end{tabular}
%* Please fill out the table in Excel using the Abrreviations specified on this page, then import the table in https://www.latex-tables.com/
% Press Auto-Booktabs, then press generate text.
% Lastly, please add the command \midpointrule after the first line, e.g.,
% Section & Tool Used & Reasoning  \\\midrule 
\end{itemize}


 \vspace{2cm}

%\noindent\rule[1ex]{30ex}{1pt} \hspace{15ex} \rule[1ex]{30ex}{1pt}\vspace{-0.2cm}\\
%\phantom{}\hspace{6ex} Ort, Datum  \hspace{37.5ex} Unterschrift \vspace{0.1cm}
%
%\newpage
%
%\begin{center} {\LARGE Erklärung} \end{center}
%
%Ich erkläre mich hiermit einverstanden, dass die vorliegende Arbeit in der Lehrstuhlbibliothek aufbewahrt wird und kopiert werden darf und habe keine Einwände gegen eine Veröffentlichung der gesamten Arbeit oder von einzelnen Teilen daraus.\vspace{0.5cm} \\
%
%Ich bin damit einverstanden, dass die RWTH Aachen folgende Daten zu meiner Person im Internet veröffentlicht: Name, Vorname, Titel der Arbeit. \\ 
%
%Mir ist bekannt, 
%\begin{itemize}
%	\item[$\Box$] 
%	dass ich diese Einwilligung jederzeit schriftlich mit Wirkung für die Zukunft widerrufen kann und meine elektronisch gespeicherten Daten unverzüglich gelöscht werden müssen und 
%	\item[$\Box$] 
%	die Aachener Verfahrenstechnik bei Widerruf meiner Einwilligung verpflichtet ist, mein Werk aus der Bibliothek des Instituts/Lehrstuhls zu entfernen.
%\end{itemize}
%
% \vspace{2cm}
%
%\noindent\rule[1ex]{30ex}{1pt} \hspace{15ex} \rule[1ex]{30ex}{1pt}\vspace{-0.2cm}\\
%\phantom{}\hspace{6ex} Ort, Datum  \hspace{37.5ex} Unterschrift \vspace{0.1cm}
\end{otherlanguage}
%
%\cleardoublepage
    \tableofcontents

% Thesis content - Add your chapters here
    \mainmatter
    \onehalfspacing
    
%   %!TEX root = ../main.tex
\chapter {Wegweiser zu den wichtigsten Latexstrukturen}\label{text:main}

\section {Tabellen}

Damit man sich die Einbettung der Tabelle im Text besser vorstellen kann, steht hier ein wenig Text. Eine Tabelle könnte folgendermaßen aussehen:

	\begin{table}[h]
		\begin{center} %Zum zentrieren der Tabelle, sieht meistens besser aus, weil die Tabellenüberschrift immer zentriert wird...
			\caption[Kurzfassung f. Inhaltsverz.]{So könnte eine Beispieltabelle aussehen} \vspace{2ex} %Die Caption ist die Tabellenunterschrift
			
			\begin{tabular}{lll}        % Die Tabelle soll drei linksbündige Spalten enthalten, 
										% alternativ r für rechtsbündig und c für zentriert, 
										% ein | fügt einen vertikalen Trennstich ein
	        \hline                      %  horizontale linie
			\textbf{Fahrzeugtyp} & \textbf{Hersteller}& \textbf{Modell}\\
			\hline                      % horizontalen Linien 
			  &\\
			  Sportwagen & Porsche & 911 GT2 \\
			   & Ferrari & Enzo \\
			  Kleinwagen & Volkswagen & Polo \\
			   & & Lupo \\
			  Familienwagen & Ford & Mondeo \\
			   & BMW & 525d \\
			  \hline
			 \end{tabular}
			\label{tbl:Autos}		% Das Label dient zum Referenzieren im Dokument
		\end{center}
	\end{table}

Unter der Tabelle geht es dann weiter mit dem Text. Es gibt auch noch andere Möglichkeiten für Tabellen. Zum Bespiel kann man Zellen zusammenfassen:

Komplexere Beispieltabelle unter Benutzung des array-Packages:


Eine Tabelle mit dem booktabs-style (Die zwei folgenden Beispiel stammen aus dem Internet\linebreak (\url{http://www.unix-ag.uni-kl.de/~fischer/blog/20070411_Tabellen_in_LaTeX/}), ich hatte keine Lust mir einen neuen Text auszudenken und die Lösung für das Multirow-Problem wollte ich auch nicht neu erfinden)\\[1.5ex] %Zeilenumbruch mit zusätzlichem Abstand (1.5 mal die Höhe des Buchstaben x)
Dann sieht die Tabelle wie folgt aus:
	
	\begin{table}[H] % Achtung: Mit diesem Befehl erzeugt man ein float-Objekt, also etwas, was von LaTeX automatisch dort platziert wird, wo es am besten aussieht, 
								% das ist nicht immer der Ort, wo man sie auch haben möchte. Abhilfe schafft z.B. das Package float und die Platzierungsoption "H". Also: \begin{table}[H]
		\caption[Schöne Tabelle]	{Das ist eine schöne Tabelle. Damit die Tabelle, wie allgemein üblich eine Überschrift und keine Unterschrift bekommt, 
															 muss der \emph{caption}-Befehl oberhalb der Tabellenbefehle, aber innerhalb der \emph{table}-Umgebung stehen. 
															 Schreibt man den \emph{caption}-Befehl unterhalb der \emph{tabular}-Umgebung, so erscheint auch einen \underline{Unter}schrift 
															 (wie bei Abbildungen üblich)}
		\begin{center} % Achtung: Dieser Befehl erzeugt zusätzlich zu dem zentrieren auch noch einen zusätzlichen vertikalen Abstand.
			\begin{tabular}{lll}
				\toprule
					Author & Title & Year 
				\\ \midrule
					\multirow{3}{*}{Philip K. Dick} & Minority Report & 1956 
				\\ \cmidrule{2-3}
		 			& Do Androids Dream of Electric Sheep? & 1968 
				\\ \cmidrule{2-3}
		 			& A Scanner Darkly & 1977 
				\\ \midrule
					\multirow{3}{*}{Neal Stephenson} & Snow Crash & 1992 
				\\ \cmidrule{2-3}
		 			& The Diamond Age & 1995
				\\ \cmidrule{2-3}
		 			& Cryptonomicon & 1999
				\\ \bottomrule\\
			\end{tabular}
		\end{center}
	\end{table} 
	
	Tabellen bei denen sich in einigen Spalten Text und in anderen Spalten Matheformeln befinden, definiert man am besten unter Benutzung des array-Paketes. Hier lässt sich am Anfang bzw. Ende einer Zelle ein beliebiger Text einfügen. das funktioniert mit ,,>\{Text\}'' bzw. ,,<\{Text\}''. Statt einem Text kann man nun einfach jeweils ein \$ einfügen und hat so eine Mathe-Spalte. Wenn man immer zwischen zwei Spalten einen bestimmten Text wie z.B. ein ,,='' haben will, kann man das durch ,,\@\{=\} erhalten. Verschönerungen wie zusätzliche Abstände etc. kann man ebenfalls einplanen. 
	
	Hier noch ein einfaches Beispiel:
	
\begin{center}
	{\centering % der Centering-Befehl erzeugt keinen zusätzlichen vertikalen Abstand.  Der Befehl und die zu zentrierenden Objekte müssen in geschweifte Klammern...
		\begin{tabular}{l>{$}c<{$} @{\;$=$\;}>{$}c<{$}}
		\toprule
		Formel 1 & \alpha & 7\cdot \beta
		\\ \midrule
		Formel 2 & x^2 & \alpha + \beta
		\\
		& & 8\cdot \alpha
		\\ \bottomrule
		\end{tabular}
	}
\end{center}

\textbf{HINWEIS:} Ein sehr hilfreiches Tool zum Erstellen von Exceltabellen ist das Excel-Plugin \glqq Excel2Latex\grqq. Damit lassen sich Tabellen sehr einfach von Excel in den Latex-Code transferieren.


\section{Mathematische Formeln}

Einzelne Formeln werden über eine \emph{Equation}umgebung
eingebunden

\begin{equation}\label{eqn:Testgleichung}
    \alpha = \beta \cdot \gamma + \delta
\end{equation}

% % Achtung: Die Eqnarray-Umgebung ist veraltet und sollte nicht mehr verwendet werden!
%Während mehrere Formeln über eine \emph{eqnarray}-Umgebung 
%übersichtlich ausgerichtet werden können.
%\begin{eqnarray}
%% \nonumber to remove numbering (before each equation)
%  \delta &=& \alpha - \beta \cdot \gamma \nonumber \\
%  \gamma &=& \frac{\alpha - \delta}{\beta}\label{eqn:Testgleichungsarray}
%\end{eqnarray}
Mehrere Formeln können mit der \emph{align}-Umgebung ausgerichtet werden:
	\begin{align}
		\delta &= \alpha - \beta \cdot \gamma & d &= f+q
		\nonumber 
			% use \nonumber or \notag to remove numbering; 
			% use \tag{OwnTag} to create a personal tag or \tag*{} to create tag without parantheses (works also in nonumbered environment)
	  \\
	  \gamma &= \frac{\alpha - \delta}{\beta}
	  \label{eqn:Testgleichungsarray}
	\end{align}

Die Nummerierung in den \emph{equation}- und \emph{align}-Umgebungen kann durch \emph{equation*} bzw. \emph{align*} unterdrückt werden.

In Gleichung~\ref{eqn:Testgleichungsarray} kann man sehen, wie man mit der \emph{align}-Umgebung Formeln anordnen kann. Grundsätzlich werden, wie bei
Abbildungen und Tabellen auch Formeln durchnummeriert. Wenn das nicht gewünscht ist, also z.B. im obigen Gleichungsarray, die
zweite Gleichung nur eine Umformung der ersten ist, reicht ein \emph{$\backslash$nonumber} vor dem Umbruch \emph{$\backslash\backslash$} um eine Nummerierung zu
unterdrücken. Wenn man die gesamte Nummerierung unterdrücken möchte, fügt man dem Befehl einfach ein * hinzu, man verwendet also z.B. $\backslash$begin$\{$align*$\}$.
Abgesetzte Formeln ohne Nummerierung können auch mit den Befehlen \textbackslash [ und \textbackslash ] erzeugt werden. 

Selbstverständlich können auch im laufenden Text Formeln wie
$\cos^{2}{x}+\sin^{2}{x}=1$ eingebunden werden.

Einheiten werden in Mathematikumgebungen mit
\emph{$\backslash$mathrm} im Gegensatz zur restlichen Formel und
Formelzeichen \textbf{nicht} kursiv dargestellt:

\begin{eqnarray}
% \nonumber to remove numbering (before each equation)
  l * b &=& A \\
  5 \mathrm{m} \cdot 2 \mathrm{m} &=& 10 \mathrm{m^{2}}
\end{eqnarray}

\section{Chemische Formeln}

	Um chemische Formeln sinnvoll zu setzen gibt es das Paket \emph{mhchem}. Dieses Paket stellt die für chemische Reaktionen nötigen Reaktionspfeile zur Verfügung 
	und setzt chemische Formeln in nicht kursiven Buchstaben.
	
	Das mhchem-Paket kann chemische Formeln wie dieses hier \ce{|\nu_A|\, A + |\nu_B|\, B <=>[oben][unten] |\nu_P|\, P + |\nu_Q|\, Q} in den Text einbinden. 
	Die Gleichungen werden am Zeilenende umgebrochen, das ist nicht immer schön. 
	
	Außerdem kann man die Befehle auch im Mathemodus einbinden und so die Gleichungen mit durchnummerieren. 
		\begin{equation}
			\ce{|\nu_A|\, A + |\nu_B|\, B <=>[oben][unten] |\nu_P|\, P + |\nu_Q|\, Q}
		\end{equation}
	
	Weitere Möglichkeiten zu dem Package findet man in der ausführlichen Dokumentation.
	
	Wenn man organische Moleküle in ihrer Struktur zeichnen will, bietet sich das Package \emph{chemfig} an. Die etwas komplizierte Benennung der Winkel entnimmt man am besten wieder der Dokumentation. Im Folgenden zwei Beispiel:
	\subsection{Trifluoroessigsäure (TFA)}
		
			Formel:\quad \ce{C2HF3O2}
			\hspace*{3em}
			\chemfig{C(-[2]F)(-[4]F)(-[6]F)-C(-[1]OH)=[7]O}
			\hspace*{3em}
			\chemfig{(-[3]F)(<:[:-150]F)(<[:-100]F)-(-[1]OH)=[7]O}\\[2ex]
			
		\subsection{Dichlormethan (DCM)}
			
			Formel:\quad \ce{CH2Cl2} 
			\hspace*{5em}
			\chemfig{C(-[5]Cl)(-[2]Cl)(<[:-70]H)(<:[:-20]H)}\\[2ex]
			

\section{Bilder}

Bilder werden mit Hilfe der ''figure''-Umgebung eingebunden. Dabei
kann die Breite festgelegt werden. Es empfiehlt sich, hier keine
absoluten Werte in Zentimetern anzugeben, sondern Bezug auf die
Textbreite zu nehmen. Das folgende Bild wurde mit 30 Prozent der
Textbreite eingebunden, \emph{0.3$\backslash$textwidth}. Alternativ kann man die Größenangeabe weglassen und das Bild in Originalgröße einbinden, die Texthöhe (textheight) oder cm-Angaben verwenden.


	\begin{figure}[h!bt]	% Requires \usepackage{graphicx}
		\begin{center}
		  \includegraphics[width=\textwidth]{rwth_avt_rgb}
		  \caption{Testbild}\label{fig:Testbild}
		\end{center}
	\end{figure}

 \begin{figure}[H]
    \centering
        \begin{subfigure}[t]{0.45\textwidth}
            \centering
            \includegraphics[width=\textwidth]{4_images/Beispiel-Graph.png}
            \caption{}
            \label{subfig:graph1}
        \end{subfigure}
        \hfill
        \begin{subfigure}[t]{0.45\textwidth}
            \centering
            \includegraphics[width=\textwidth]{4_images/Beispiel-Graph.png}
            \caption{}
            \label{subfig:graph2}
        \end{subfigure}
    \caption{Graphen für Ergebnisse werden mit Origin erstellt. Ein Tutorial kann unter https://cvtwiki.avt.rwth-aachen.de/wiki/doku.php?id=sop:data\_management:origin gefunden werden. Achtung, der Link lässt sich nur aus der pdf öffnen, wegen des \_.}
    \label{fig:graphs}
\end{figure}


Mögliche Bildformate sind *.jpg, *.png, *.gif, *.pdf.

Das Einbinden von Bildern erfolgt mit der  $\backslash begin\{ figure\}...\backslash end\{ figure\}$-Umgebung als sogenanntes Float-Element, mit dem die Position von Latex selber angepasst werden kann. Die Option [H] zwingt das Float-Element genau an die stelle im Text an der das Element im Quelltext eingebunden wurde.

\section{Fußnoten}
Ein kurzer Beispieltext mit Fußnoten:

Grundsätzlich gilt: Ein Blick in die Literatur spart viel mehr
Zeit als sie kostet! Und damit sie noch weniger Zeit kostet, sind
die wichtigsten Wege zur Recherche und Beschaffung von Literatur
an der RWTH im AVT-Wiki\footnote{Wikis haben den Vorteil das sie
kontinuierlich verbessert werden können, während dieses Dokument
doch eher statischer Natur ist} zusammengestellt. Die angegeben
Informationsquellen stellen nur eine Übersicht über die
Möglichkeiten dar und haben nicht den Anspruch auf
Vollständigkeit, können aber mit Hilfe aller Nutzer diesen Zustand
anstreben. Zu finden ist der Wegweiser zur Literaturbeschaffung
auf
\url{https://wiki.avt.rwth-aachen.de}
zu finden.\footnote{Diese URL ist nur an Institutsrechnern oder
über den entsprechenden Login erreichbar}

\section{Graphen und Zeichnungen}

Graphen können ohne viel Aufwand mit PGFplots erstellt werden. Beispiele gibt es unter \url{http://pgfplots.sourceforge.net/gallery.html}.
\begin{center}
\begin{tikzpicture}
	\begin{axis}[view/h=-30]
	\addplot3[
		surf,
		shader=flat,
		samples=30,
		domain=0:2,y domain=0:2] 
		{sin(deg(x^2+y^2))};
	\end{axis}
\end{tikzpicture}
\end{center}

Zeichnungen können unkompliziert mit Tikz angefertigt werden. Siehe \url{http://www.texample.net/tikz/examples/}.
\begin{center}
\begin{tikzpicture}[->,>=stealth',shorten >=1pt,auto,node distance=4cm,semithick]
  \node[state] (A){$s_0$};
  \node[state] (B) [above right of=A] {$s_1$};
  \node[state] (C) [below right of=B] {$s_2$};

  \path 
    (A) edge [bend left] node{0,2}(B)
        edge node {0,2}(C)
    edge [loop left] node{0,6}(A)
    (B) edge [loop above] node{0,5}(B)
        edge [bend left] node{0,4}(C)
    edge node{0,1}(A)
    (C) edge node{0,2}(B)
        edge[bend left] node{0,3}(A)
    edge[loop right] node{0,5}(C);
\end{tikzpicture}
\end{center}


\section{Internetadressen}

Abhilfe bei (fast) allen Problemen bietet der Befehl \emph{url}:

\url{https://avt.rwth-aachen.de}

Dann sollte ein Zeilenumbruch erfolgen. Achtung, auch hier gibt es ein Sonderzeichen, das Probleme machen kann und das ist das \%, da man damit in \LaTeX\ den Rest der Zeile auskommentiert. Hier muss man wieder \textbackslash \% schreiben.

\section {Listen}
Hier sei zunächst auf einige nützliche Zusammenstellungen von
Stoffdaten verwiesen, um das Einfügen einer unnummerierten Liste
zu demonstrieren:
\begin{itemize}
    \item Eine sehr gute Übersicht findet man auf den Seiten des \emph{\textbf{National Institute of Standards and Technology}} (www.nist.gov) unter Information for Researchers.
    \item \emph{\textbf{Perry's Chemical Engineers' Handbook}}, Kapitel 2 (IVT-Bibliothek)
    \item \emph{\textbf{CRC Handbook of Chemistry and Physics}} (IVT-Bibliothek)
    \item \emph{\textbf{Handbook of Electrolyte Solutions}} (zweibändig, unter Signatur E7146 im Großen Lesesaal der HSB)
    \item \emph{\textbf{Gmehling-Onken}}, Sammlung Thermochemischer Stoffdaten, vor allem Gas-Flüssig Gleichgewichte (mehrbändig, unter Signatur E0218 im Großen Lesesaal der HSB sowie in der Bibliothek der AVT.TVT)
    \item In der Chemiebibliothek in der Prof.-Pirlet-Str. finden sich weitere Material- und Datensammlungen. Das Stöbern lohnt sich.
    \item Im Programmpaket \emph{\textbf{ASPEN Plus}} ist eine umfangreiche Stoffdatenbank integriert, die auch unabhängig von einer Prozesssimulation abgefragt werden kann.
    \item   \begin{enumerate}
        \item \textbf{Selbstverständlich können Listen auch geschachtelt
        werden}
        \item Nummerierte Listen sind auch möglich, nur heißt dann die Umgebung nicht ''itemize'' sondern ''enumerate''
            \end{enumerate}

\end{itemize}

		\section {Gliederung}
			Wichtig ist nicht nur die technisch richtige Gliederung in Tex,
			sondern auch das richtige bezeichnen in der Arbeit. So heißt
			dieses Kapitel ''Wegweiser zu den wichtigsten Latexstrukturen''
			und sollte auch so benannt werden.
			
			Ist das Unterkapitel ''Gliederung'' gemeint, sollte es bei
			Verweisen auch so genannt werden. Gleiches gilt für
			\subsection {Abschnitte}
			Abschnitte
			\subsubsection {Unterabschnitte}
			Unterabschnitte und
			\paragraph {Absätze}
			Absätze
			
			Diese Gliederungseinheiten können übrigens auch im Code mit einem
			* hinter dem Namen versehen werden. Dann wird eine Nummerierung
			unterdrückt und die Überschrift nicht im Inhaltsverzeichnis
			gelistet. Sinn macht das aber nur bei Kapiteln, Unterkapiteln und
			Abschnitten, da nur diese nummeriert werden.

		\section {Umgebungen}
			Tabellen, Bilder und Gleichungen werden in Latex in der Regel in
			Umgebungen geschrieben. Bei Tabellen und Bildern enthalten diese
			oft den \textbf{$\backslash$caption$[x]$\{y\}}-Befehl. y steht in
			diesem Fall für die Bild-/Tabellen Unter- bzw. Überschrift,
			abhängig davon, ob der Befehl in der Umgebung vor oder nach dem
			Objekt
			steht.\\
			x ist optional und steht für eine Kurzbeschreibung, falls y zu
			lang für das Abbildungs- oder Tabellenverzeichnis ist.
			
			\subsection {Referenzen}
			Um auf bestimmte Textstellen oder Objekte zu verweisen, bedient
			man sich in Latex sogenannter Labels und Referenzen. Ein Label
			wird mit dem Befehl \textbf{$\backslash$label\{name\}} gesetzt.
			Der Name sollte hier möglichst intuitiv gesetzt werden. Es ist
			sinnvoll hier auch den Typ des verwiesenen Objektes einzubeziehen,
			z.B. ist es empfehlenswert den Namen bei Gleichungen mit ''eqn:''
			zu beginnen, bei Bildern mit ''fig:'' und bei Tabellen mit
			''tab:''.
			
			Sehr wichtig ist, dass bei Bildern und Tabellen die
			\textbackslash label-Anweisung unmittelbar auf die
			\textbackslash caption-Anweisung und bei Gleichungen noch in derselben Zeile steht, da sonst auf die Textstelle statt auf das
			Objekt verwiesen wird.
			
			Neben den nummerierten Referenzen mit \textbf{$\backslash$ref\{name\}}, wie z.B. auf Gleichung
			\ref{eqn:Testgleichungsarray}, ist mit \textbf{$\backslash$pageref\{name\}} der Hinweis möglich, dass
			besagte Gleichung auf Seite \pageref{eqn:Testgleichungsarray} zu finden ist. Besonders interessant ist 
			der Befehl \textbf{\textbackslash eqref\{name\}} für mathematische Gleichungen, da hier automatisch die Klammern um die
			Referenznummer mitgesetzt werden.
			
		
		\section {Hinweise zum Literaturverzeichnis}
			
			Die für das Literaturverzeichnis verwendeten Daten, werden in einem externen file mit der Endung .bib gespeichert. Diese Datei kann theoretisch in jedem Texteditor erstellt und bearbeitet werden, solange man sich an die vorgegebene Syntax hält. 
			Einfacher ist die Verwendung des Open Source Tool \emph{JabRef}, das eine Maske mit Feldern bereitstellt, bei denen auch vorgegeben ist, welche Angaben benötigt werden und welche nicht unbedingt. Außerdem kann man in JabRef auch Zusatzinformationen wie den Abstract der Paper speichern und die Literatur verwalten.
			
			Das zitieren erfolgt über den sogenannten Bibtex-Key mit dem Befehl \textbackslash cite\{BibtexKey\}. Das Literaturverzeichnis wird automatisch nach \textbf{zweimaligem} kompilieren aktuell erstellt (Beim ersten Kompilieren werden die Daten der .bib-Datei ausgelesen und beim zweiten Kompilieren wirklich verknüpft). 
			
			
			Und so sieht das dann praktisch aus: Das hier ist eine sehr interessante Quelle \cite{Abel2000}, man sollte aber auch dieses Paper unbedingt gelesen haben \cite{Marquardt2005}. Dabei ist die autovervollständigen Funktion vom TeXnicCenter durchaus hilfreich. Wenn man einen der Befehle \textbackslash cite oder \textbackslash ref verwendet, bekommt man nach dem Tippen der ersten Zeichen des Keys Vorschläge, die man mit Strg + Leertaste annehmen kann.


		    
%  %authors: Arne Lüken and Lukas Griesberg 02/2022
%please send comments and additions to Arne Lüken arne.lueken@avt.rwth-aachen.de
\chapter{Wissenschaftliches Schreiben}

\section{Schreibstil}
    Ein wissenschaftlicher Text sollte in der Lage sein, komplexe Themen verständlich zu vermitteln. Sätze und Formulierungen sind hierfür möglichst kurz und klar zu gestalten. Vermeiden Sie unnötige Informationen und Wörter. Alle Wörter, die ausgelassen werden können, ohne den Sinn zu verändern, sollten ausgelassen werden. Der Sprachstil ist sachlich zu halten; ein Spannungsaufbau sollte nicht stattfinden. Verwenden Sie nach Möglichkeit Aktiv anstatt Passiv. Speziell im Englischen macht es den Text lebendiger:
    \begin{compactitem}
        \item \underline{Passiv:} "Der Katalysator ist auf der Nickelelektrode aufgebracht."
        \item \underline{Aktiv:} "Die Elektrode besteht aus einer mit Katalysator beschichteten Nickelplatte."
    \end{compactitem}
    Verwenden Sie das Präsens für allgemeine Informationen, Bezug auf Graphen oder Abbildungen, sowie bei der Interpretation von Ergebnissen und das Präteritum für das Zitieren von Forschungsergebnissen und bei Bezug auf eigene Messungen und Ergebnisse (z.B. "...wurde gemessen."). Achten Sie zudem auf eine einheitliche Wortwahl. Sprechen Sie beispielsweise nicht abwechselnd von "Mikrogelen", "Hydrogelen" oder "Partikeln", sondern beschränken Sie sich auf einen dieser Begriff.

\section{Struktur von Text und Graphen}

    Bevor mit dem Schreiben begonnen wird, sollte ein roter Faden erstellt werden. Zunächst sollten die Kapitel und Unterkapitel grob festgelegt werden, um einen Überblick über die Struktur der Arbeit zu erhalten. Falls Sie Unterkapitel verwenden, müssen es jeweils mindestens zwei Unterkapitel sein (z.B. 1.3.1. und 1.3.2.). Meist bietet es sich an, vor mehreren Unterkapiteln eine kurze Einleitung des Oberkapitels zu setzen. In der Regel sollten Unter-Unter-Kapitel (1.3.1.) ausreichend sein, um die Übersichtlichkeit der Arbeit zu gewährleisten.
    
    Bevor Sie zu schreiben anfangen, bietet es sich an, die Abbildungen und Graphen, die sie verwenden wollen, einzufügen und als Ausgangspunkt für Ihren Text zu nehmen. Sollten diese noch nicht fertig sein, kann es helfen, die erwarteten Ergebnisse als Skizze/Handzeichnung einzufügen und am Ende durch die fertigen Abbildungen zu ersetzen. Wichtig ist, dass jede verwendete Abbildung im Text referenziert wird. Ergebnisse, die Sie für wichtig erachten, jedoch nicht direkt ansprechen oder interpretieren, gehören in den Appendix. Alle relevanten Informationen sollten im Graphen oder der Bildunterschrift, welche bei Papern über mehrere Zeilen gehen darf, enthalten sein. Der Leser sollte ohne den Fließtext zu lesen, erkennen, was er auf der Abbildung sieht und was diese ausmacht, wie etwa relevante Versuchsbedingungen.
    
    Bei der Erstellung von Abbildungen und Graphen ist es empfehlenswert, Farbe nur falls unbedingt nötig zu verwenden. Es sollten keine Informationen verloren gehen, falls die Arbeit in schwarz-weiß gedruckt wird. Abbildungen sollten möglichst kontrastreich gestaltet werden. Falls farbige Abbildungen unvermeidbar sind, sollten Rot und Grün möglichst vermieden werden, da diese für farbenblinde Menschen problematisch sein können. Für Beschriftungen innerhalb von Abbildungen verwenden Sie bitte kurvige Linien (siehe Abbildung~\ref{fig:Beispielgraph}). Graphen sollten mit der Software \href{https://www.originlab.com/}{Origin} erstellt werden. Für Origin gibt es Vorlagen und eine \href{https://cvtwiki.avt.rwth-aachen.de/wiki/doku.php?id=sop:data_management:origin}{\textbf{\textit{Origin-Anleitung}}} im CVTwiki. Bilder und Schemata sollten immer einen Maßstabsbalken enthalten.
    
    \begin{figure}[h]	% Requires \usepackage{graphicx}
    		\begin{center}
    		  \includegraphics[width=\textwidth]{4_images/Beispiel-Graph.png}
    		  \caption[Beispielgraph]{Ein Beispielgraph mit Gestaltungshinweisen, die in \href{https://cvtwiki.avt.rwth-aachen.de/wiki/doku.php?id=sop:data_management:origin}{Origin} einfach und direkt angepasst werden können. Das Origintutorial kann im CVT Wiki unter https://cvtwiki.avt.rwth-aachen.de/wiki/doku.php?id=sop:data\_management:origin gefunden werden. Achtung, dieser Link funktioniert nicht im Code Editor wegen \_.}\label{fig:Beispielgraph}
    		\end{center}
    \end{figure}
    
    Um eine logische, nachvollziehbare Textstruktur zu erhalten, bietet sich die Verwendung von Absätzen an. Hierbei sollte jeder Absatz jeweils eine Idee behandeln und diese ausführen. In der Regel hat ein Absatz eine Länge von 5~-~10 Zeilen, so dass 3~-~5 Absätze auf eine Seite passen. Jeder Absatz sollte daher grob wie folgt strukturiert sein:
    \begin{compactitem}
        \item \underline{Topic Sentence:} Einleitender Satz, der die Grundaussage des Absatzes darlegt. Er beinhaltet sowohl das Thema (z.B. Mikrogele) als auch die Leitidee (z.B. anisometrisches Schwellen).
        \item \underline{Supporting Sentences:} Erklärung oder Unterstützung der Hauptidee durch Details in mehreren Sätzen.
        \item \underline{Concluding Sentence:} Fasst die wichtigsten Punkte zusammen und signalisiert das Ende des Absatzes. Er ist ähnlich strukturiert wie der Topic Sentence, beinhaltet aber mehr Informationen.
    \end{compactitem}
    
    
\FloatBarrier
\section{Gliederung}
Wissenschaftliche Arbeiten gliedern sich in der Regel in folgende Abschnitte:
\begin{enumerate}
    \item[] \textbf{Abstract} \vspace{-4 mm}
    \begin{compactitem}
        \item Sehr kurze Zusammenfassung der Thesis; Überblick für potenzielle Leser
        \item Inhalt:
        \begin{compactitem}
            \item Warum untersucht die Arbeit das Thema? (Motivation)
            \item Wie wurde das Thema behandelt? (Methoden)
            \item Was sind die Ergebnisse? (Ein bis zwei der wichtigsten Erkenntnisse, die am Ende herauskommen)
            %ONE general sentence
            %ONE sentence about the problem
            %state the objective
            %state the methods
            %state most significant results
            %take-home message
        \end{compactitem}
        \item Struktur:
        \begin{compactitem}
            \item Maximal 200 Wörter.
            \item Gegenwart, unpersönlicher Stil (vermeiden Sie \textit{ich} oder \textit{wir}).
        \end{compactitem}
    \end{compactitem}
    \item \textbf{Introduction / Motivation} \vspace{-4 mm}
    \begin{compactitem}
        \item Inhalt:
        \begin{compactitem}
            \item Warum wird das Thema untersucht? Was ist der Hintergrund der Arbeit? (Keine langen Ausschweifungen; kurz und prägnant Motivation erklären)
            \item Wer sind die Interessensgruppen? (z.B. für welche Anwendungen ist es interessant? Gibt es Produkte, die mithilfe der Ergebnisse entwickelt werden könnten? etc.)
            \item Um welche Problemstellung geht es genau und wie wird diese angegangen? 
            \item Was ist die Stuktur der Arbeit?
        \end{compactitem}
        \item Struktur:
        \begin{compactitem}
            \item Ca. eine Seite, höchstens zwei.
            \item Keine Unterkapitel.
        \end{compactitem}
    \end{compactitem}
    \item \textbf{Theoretical Background and State of the Art} \vspace{-4 mm}
    \begin{compactitem}
        \item Liefert Grundlage für die spätere Diskussion und Interpretation der Ergebnisse
        \item Inhalt:
        \begin{compactitem}
            \item Was sind die Grundlagen und Hintergründe des behandelten Themas?
            \item Wie ist der bisherige Stand der Forschung?
            \item Was wurde im behandelten Thema bereits am Institut getan?
            \item Was sind Schnittmengen zu anderen Forschungsfeldern? Wie werden die Methoden sonst angewendet?
        \end{compactitem}
        \item Struktur:
        \begin{compactitem}
            \item Unterkapitel ergeben Sinn.
            \item Es sollte eine visuelle Darstellung des State-of-the-art erstellt werden (z.B. chronologisch wie in Abbildung \ref{fig:ChronologicStateOfTheArt}).
        \end{compactitem}
    \end{compactitem}
    \item \textbf{Concept / Hypothesis} (optional) \vspace{-4 mm}
    \begin{compactitem}
        \item Inhalt:
        \begin{compactitem}
            \item Im Gegensatz zur Einleitung sollen hier aufbauend auf der wissenschaftlichen Lücke des State of the Arts die der Arbeit zugrundeliegenden Hypothesen detailliert dargestellt werden. Darauf aufbauend werden dann die spezifischen angewandten Methoden und die erwarteten Ergebnisse aufgelistet. Was? Warum?
        \end{compactitem}
        \item Struktur:
        \begin{compactitem}
            \item Circa eine Seite.
        \end{compactitem}
    \end{compactitem}
    \item \textbf{Materials and Methods} \vspace{-4 mm}
    \begin{compactitem}
        \item Genaue Beschreibung des verwendeten Versuchsaufbaus und -durchführung, damit die Versuche \textit{reproduzierbar} sind.
        \item Inhalt:
        \begin{compactitem}
            \item Kapitel sollte enthalten:
                 Materialien (inkl. Hersteller), Verwendete Software,
                 Mengen / Konzentrationen,
                 Versuchsaufbau,
                 Versuchsvorbereitung,
                 Kalibrierung,
                 Versuchsprotokolle
            \item \underline{Der Versuch sollte auf Basis dieses Kapitels replizierbar sein}
            \item Gehen Sie auf eventuell vorgenommene Veränderungen am Versuchsaufbau ein, die stattgefunden haben, und nennen Sie Probleme, die aufgetreten sind.
            \item \underline{Beachte:} Keine Ergebnisse und keine Bewertung in diesem Kapitel.
        \end{compactitem}
        \item Struktur:
        \begin{compactitem}
            \item Unterkapitel sind möglich.
            \item Tabellen können hilfreich sein.
            \item Versuchsaufbau als beschriftetes Foto / Skizze / Verfahrensfließbild.
            \item Kurz und prägnant.
        \end{compactitem}
    \end{compactitem}
    \item \textbf{Results and Discussion} \vspace{-4 mm}
    \begin{compactitem}
        \item Inhalt:
        \begin{compactitem}
            \item Alle relevanten Ergebnisse sollten aufgeführt werden und in den Kontext der Arbeit gebracht werden. Nicht für das Ergebnis der Arbeit relevante Ergebnisse sollten im Anhang untergebracht werden.
            \item Verwendung von Graphen und Schemata.
            \item Ergebnisse werden zunächst beschrieben (Results) und anschließend im wissenschaftlichen Kontext (Bezug zu State-of-the-art) analysiert und interpretiert (Discussion).
        \end{compactitem}
        \item Struktur:
        \begin{compactitem}
            \item Die Ergebnisse sollen nicht chronologisch, sondern aufeinander aufbauend dargestellt werden, wie man auch eine Geschichte erzählen würde.
            \item Der rote Faden sollte anhand von Abbildungen (Graphen) erkennbar sein.
            \item Results und Discussion können als getrennte Kapitel behandelt oder zusammengefasst werden.
        \end{compactitem}
    \end{compactitem}
    \item \textbf{Conclusion and Outlook} \vspace{-4 mm}
    \begin{compactitem}
        \item Zusammenfassung der zentralen Ergebnisse für jemanden, der die Arbeit nicht komplett gelesen hat.
        \item Inhalt:
        \begin{compactitem}
            \item Ergebnisse werden erneut aufgeführt und die in "Results and Discussion" erreichten Schlüsse werden wiederholt. 
            \item Es werden keine neuen Diskussionen / Punkte erörtert.
            \item \underline{Outlook:} Mögliche Ansätze für weiterführende Forschung. Vorschläge und Ideen, wie es weitergehen kann.
        \end{compactitem}
        \item Struktur:
        \begin{compactitem}
            \item 1-2 Seiten; je maximal eine Seite für Conclusion und Outlook.
            \item Kann je nach Bedarf in einzelne Kapitel für Conclusion und Outlook aufgeteilt werden.
        \end{compactitem}
    \end{compactitem}
\end{enumerate}

\begin{figure}[h!bt]	% Requires \usepackage{graphicx}
		\begin{center}
		  \includegraphics[width=1\textwidth]{4_images/Beispiel_StateArt.png}
		  \caption{Beispielhafte Visualisierung des State-of-the-art im Bereich "Magnetresonanztomographie (MRI) in Membranfiltration" aus den Jahren 2002 - 2018.}\label{fig:ChronologicStateOfTheArt}
		\end{center}
\end{figure}
 
    
    \chapter{Einleitung}
Hier beginnt deine Einleitung...
    \chapter{Theoretischer Hintergrund}
Beschreibe hier die Theorie...
%    \include{2_chapter/3_Concept/Hypothesis} %optional
    \chapter{Materialien und Methoden}
Beschreibe hier die Methoden...
    \chapter{Ergebnisse und Diskussion}
Präsentiere deine Ergebnisse...
    \chapter{Fazit und Ausblick}
Ziehe hier dein Fazit...
    
% Nomenclature, bibliography, etc.
%    Hier steht deine Danksagung...
    %!TEX root = ./main.tex
\tikzset{external/prefix=tikz/Nomenclature/}
% Abkürzungen definieren
\newacronym{cvd}{CVD}{Chemical Vapour Deposition}
\newacronym{lpcvd}{LPCVD}{Low Pressure \acrshort{cvd}}
\newacronym{apcvd}{APCVD}{Atmospheric-Pressure \acrshort{cvd}}
\newacronym{mocvd}{MOCVD}{Metal-organic \acrshort{cvd}}
\newacronym{ald}{ALD}{Atomic Layer Deposition}
\newacronym{cvi}{CVI}{Chemical Vapour Infiltration}
\newacronym{tgcvi}{TG-CVI}{Thermal-Gradient Chemical Vapour Infiltration}
\newacronym{sem}{SEM}{Scanning Electron Microscopy}
\newacronym{fgm}{FGM}{Functionally Graded Material}
\newacronym{edx}{EDX}{Energy Dispersive X-ray}
\newacronym{psi2}{PSI-2}{Plasma Surface Interaction - 2}
\newacronym{cul}{Cu L$_{\alpha 1,2}$}{Copper L alpha$_{1,2}$ emission line}
\newacronym{edm}{EDM}{Electrical Discharge Machining}
\newacronym{se}{SE}{Secondary Electron}
\newacronym{bse}{BSE}{Backscattered Electron}
\newacronym{il}{IL}{In Lens Electron}
\newacronym{z}{Z}{Atomic number}
\newacronym{wilma}{WILMA}{Wolfram Infiltrations Maschine}
\newacronym{ccz}{\ch{CuCrZr}}{copper chromium zirconium alloy}
\newacronym{w}{\ch{W}}{tungsten}
\newacronym{wfw}{W$_f$/W}{tungsten fibre reinforced tungsten}
\newacronym{wf}{W$_f$}{tungsten fibre}
\newacronym{iter}{ITER}{International Thermonuclear Experimental Reactor}
\newacronym{elm}{ELM}{Edge-Localised Mode}
\newacronym{naoh}{\ch{NaOH}}{sodium hydroxide}
\newacronym{naf}{\ch{NaF}}{sodium fluoride}
\newacronym{hf}{\ch{HF}}{hydrogen fluoride}
\newacronym{h}{\ch{H2}}{hydrogen}
\newacronym{lrp}{LRP}{liquid ring pump}
\newacronym{wf6}{\ch{WF6}}{tungsten hexafluoride}
\newacronym{ar}{\ch{Ar}}{Argon}
\newacronym{cvdw}{CVD-W}{\acl{cvd} coating with \acl{w}}
% --- Zusätzlich eigene Glossare für Symbole und Indizes ---


% Symbole
\newglossaryentry{theta}{type=symbols,
  name={\ensuremath{\Theta}},
  description={Contact angle between droplet and surface},
  sort=theta,
  symbol={\ensuremath{^\circ}}}

\newglossaryentry{Ea}{type=symbols,
  name={\ensuremath{E_\text{A}}},
  description={Activation energy},
  sort=Ea,
  symbol={J}}

% Indizes
\newglossaryentry{max}{type=indices,
  name={max},
  description={Maximum value}}

\newglossaryentry{f}{type=indices,
  name={f},
  description={Fluorescence}}


    % Verzeichnisse
    % Abkürzungsverzeichnis ins TOC

    \printglossary[type=\acronymtype, title=List of Abbreviations, toctitle=List of Abbreviations]
    \printglossary[type=symbols, title=List of Symbols, toctitle=List of Symbols]
    \printglossary[type=indices, title=List of Indices, toctitle=List of Indices]

    \printbibliography[heading=bibintoc]
    \listoffigures
    \listoftables
    \appendix
    \chapter{Anhang}
Hier steht dein Anhang...
    
    
\end{document}
